\documentclass[dvipdfmx,twoside]{jsarticle}
\usepackage{amsmath,amssymb}
\usepackage{CJKutf8}
\usepackage{okumacro}
\usepackage{booktabs}
\usepackage{array}
\usepackage{xcolor}
\usepackage{colortbl}
\usepackage{tipa}
\usepackage{tikz}
\usepackage{fancybox}
\usepackage{wasysym}
\usepackage{pifont}
\usepackage{graphicx}
\usepackage{wrapfig}
\usepackage{geometry}
\geometry{margin=2cm}
\usepackage{fancyhdr}
\usepackage[utf8]{inputenc}
\usepackage{CJK}
% 设置页面样式
\pagestyle{fancy}
\fancyhf{}
\renewcommand{\headrulewidth}{0pt}
\fancyhead[RO]{数学–\thepage}
\fancyhead[LE]{数学–\thepage}

% 重新定义 plain 样式
\fancypagestyle{plain}{%
  \fancyhf{}%
  \renewcommand{\headrulewidth}{0pt}%
  \fancyhead[RO]{数学-\thepage}%
  \fancyhead[LE]{数学-\thepage}%
}
% 自定义A方框命令
\newcommand{\abb}[1]{%
\begin{tikzpicture}[baseline]
\node[draw=black, 
      rectangle, 
      minimum width=0.8cm, 
      minimum height=0.3cm, 
      fill=gray!25, 
      font=\bfseries,
      line width=1pt,
      inner sep=2pt,
      anchor=base] {#1};
\end{tikzpicture}%
}
\newcommand{\ab}[1]{%
\begin{tikzpicture}[baseline]
\node[draw=black, 
      rectangle, 
      minimum width=0.8cm, 
      minimum height=0.3cm, 
      font=\bfseries,
      line width=1pt,
      inner sep=2pt,
      anchor=base] {#1};
\end{tikzpicture}%
}

\newcommand{\maru}[1]{\tikz[baseline=-0.7ex]{
    \node[shape=circle,draw,inner sep=1pt,minimum size=5pt,anchor=center] {\footnotesize #1};}}
\definecolor{headercolor}{RGB}{220,220,220}
\definecolor{rowcolor1}{RGB}{245,245,245}
\definecolor{rowcolor2}{RGB}{255,255,255}

% \title{\vspace{-1.5cm} 2025年8月羚課文科数学月考 }
% \author{\textnosfal{Linc\ -\ 伊}}
\date{}
\begin{document}
%第一部分
\begin{CJK}{UTF8}{ipxm}  % 使用ipxm字体
  %月考表纸部分
\begin{center}

\vspace*{5cm}

% 学校logo
\includegraphics[width=5cm]{pics/1.jpg}

\vspace{2cm}

% 主标题
{\fontsize{24}{30}\selectfont\bfseries\sffamily
2025年8月\\
\vspace{1em}
羚課文科数学月考
}

\end{center}
\newpage
\noindent
\begin{tikzpicture}
\node[draw, rectangle, minimum width=1cm, minimum height=0.8cm, font=\Huge] {I};
\end{tikzpicture}
\\
\\
\textbf{問1}\qquad $ x^2+y^2=4 $のとき、 $ 2x+y $の最大値は$\ab{\textsf{A}}\sqrt{\ab{\textsf{B}}}$\\

2次関数$y = x^2 + 6x + 5$ のグラフを点$(1,3)$に関して対称移動してできるグラフの方程式は\\
$$ y= \ab{\textsf{C}}\ x^2+\ab{\textsf{D}}\ x+\ab{\textsf{E}}$$
\\
\\
\\
\noindent
\textbf{問2}\qquad $3a + 1$ が $a^2 + 5$ の約数となるような自然数 $a$ を求めよう。

\vspace{2em}

$3a + 1 = b$ とする。このとき

\vspace{1em}

\begin{center}
$a^2 + 5 = \dfrac{b^2 - \ab{\textbf{F}}\ b + \ab{\textbf{GH}}}{\ab{\textbf{I}}}$ \hspace{2em} $\cdots\cdots$ \maru{1}
\end{center}

\vspace{1em}

である。また、$b$ は $a^2 + 5$ の約数であるから、$a^2 + 5$ はある自然数 $c$ を用いて

\vspace{1em}

\begin{center}
$a^2 + 5 = bc$ \hspace{2em} $\cdots\cdots$ \maru{2}
\end{center}

\vspace{1em}
と表される。\maru{1}、\maru{2} から

\vspace{1em}

\begin{center}
$b\left(\ab{\textsf{J}}\ c - b + \ab{K}\right) = \ab{LM}$
\end{center}

\vspace{1em}

を得る。したがって、$b$ は \abb{LM} の約数である。この中で、$a$ が自然数となるのは $b =$ \ab{NO}

\vspace{1em}

である。したがって、$a = \ab{PQ}$ である。
\newpage
\begin{center}
- 計算欄 (memo) -
\end{center}
\newpage
%第二部分
\noindent
\begin{tikzpicture}
\node[draw, rectangle, minimum width=1cm, minimum height=0.8cm, font=\Huge] {II};
\end{tikzpicture}
\\
\\
\noindent
\textbf{問1}\qquad 大きさの異なる4枚のカードがある。これらのカードに赤、黒、青、黄の色を塗る。ただし、\\
\hspace*{1em}どのカードにも1つの色のみを使い、また同じ色のカードが2枚以上あってもよいものとする。

\vspace{1em}

(1) \hspace{1em} 全部で \ab{\textsf{ABC}} 通りの塗り方がある。

\vspace{1em}

(2) \hspace{1em} 全部の色を使う塗り方は \ab{\textsf{DE}} 通りある。

\vspace{1em}

(3) \hspace{1em} 2枚は赤で、1枚が黒、1枚が青となるような塗り方は \ab{\textsf{FJ}} 通りある。

\vspace{1em}

(4) \hspace{1em} 3つの色を使う塗り方は \ab{\textsf{GHI}}  通りある。

\vspace{1em}

(5) \hspace{1em} 2つの色を使う塗り方は \ab{\textsf{JK}} 通りある。\\
\\
\\
\\
\textbf{問2}\qquad 2つの放物線\\
\\
$\ell: \quad y = ax^2 + 2bx + c$\\
$m: \quad y = (a+1)x^2 + 2(b+2)x + c + 3$\\
\\
\begin{wrapfigure}{r}{0.4\textwidth}
\vspace{-8em}
\begin{tikzpicture}[scale=0.8]
\draw[->] (-3,0) -- (3,0) node[right] {$x$};
\draw[->] (0,-2) -- (0,3) node[above] {$y$};
\node[below left] at (0,0) {O};
\node[above left] at (-2,0) {A};
\node[below left] at (-1.5,-1.5) {B};
\node[above right] at (-1,0) {C};
\node[below right] at (1,0) {D};
\fill (-2,0) circle (2pt);
\fill (-1.5,-1.5) circle (2pt);
\fill (-1,0) circle (2pt);
\fill (1,0) circle (2pt);
\end{tikzpicture}
\end{wrapfigure}

\noindent
を考える。点 A, B, C, D が右図のような位置関係に
\noindent
あるとする。このとき，この 2 つの放物線のうち，
\noindent
一方は，3 点 A, B, C を通り，もう一方は，3 点 B,
\noindent
C, D を通るとする。

\vspace{1em}
\noindent
(1)\quad 3 点 A, B, C を通る放物線は \ab{\textsf{L}} である。ただし、\\
\ab{\textsf{L}} には，次の\maru{0}か\maru{1}のどちらか適するものを選びなさい。

\vspace{0.5em}

\begin{center}
\begin{tabular}{cc}
\maru{0} & 放物線 $\ell$ \\
\maru{1} & 放物線 $m$
\end{tabular}
\end{center}

\vspace{1em}
\noindent
(2)\quad 2 つの放物線 $\ell, m$ は，どちらも 2 点 B, C を通るので，点 B, C の座標は，2 次方程式

\[x^2 + \ab{\textsf{M}}\ x + \ab{\textsf{N} }= 0\]

の解である。よって，点 B の $x$ 座標は \ab{\textsf{OP}}，点 C の $x$ 座標は \ab{\textsf{QR}} である。

\vspace{1em}
\noindent
(3)\quad 特に，AB = BC, CO = OD のとき，$a, b, c$ の値を求めよう。\\

\quad 2 点 C, D は $y$ 軸に関して対称であるから，$b = \ab{\textsf{S}}$ である。また，AB = BC より，\\

直線 $x = \ab{\textsf{TU}}$ が\ \abb{L}\ の軸である。したがって，$a = -\ \ab{\textsf{V}}$ である。よって，\\


$c = \ab{\textsf{W}}$ である。
\newpage
\begin{center}
- 計算欄 (memo) -
\end{center}
\newpage
%第三部分
\noindent
\begin{tikzpicture}
\node[draw, rectangle, minimum width=1cm, minimum height=0.8cm, font=\Huge] {III};
\end{tikzpicture}
\\

$n$ は 2 桁の自然数であり、$n^3$ を 66 で割ったときの余りは $n$ であるという。このような $n$
\noindent
の個数と、このような $n$ のうち素数であるものを求めよう。

\vspace{2em}

条件より、$n^3$ を 66 で割ったときの商を $p$ とすると

\vspace{1em}

\hspace{3em} $n^3 = \ab{\textsf{AB}}\ p + n$ \hspace{2em} $\left(0 < n \leq \ab{CD}\right)$

\vspace{1em}
\noindent
と表せる。これを変形して

\vspace{1em}

\hspace{3em} $n(n - 1)(n + 1) = \abb{\textsf{AB}}\ p$

\vspace{1em}
\noindent
を得る。

\vspace{1em}
ここで、$n - 1, n$ のどちらか一方は \ab{\textsf{E}} の倍数、$n - 1, n, n + 1$ のうち 1 つは \ab{\textsf{F}} の倍数であり、\\

\noindent
\abb{\textsf{E}} と \abb{\textsf{F}} は互いに素であるから、$n(n - 1)(n + 1)$ は \ab{\textsf{G}} の倍数である。\\

\noindent
ただし、$1 < \abb{\textsf{E}} < \abb{\textsf{F}} < \abb{\textsf{G}}$ とする。よって、$n - 1, n, n + 1$ のいずれか
\vspace{1em}
1 つが \ab{\textsf{HI}} の倍数となる場合を考えればよい。

\vspace{1em}

いま、$n \leq \abb{\textsf{CD}}$ であるから、$n - 1$ が \abb{\textsf{HI}} の倍数である $n$ の個数は \ab{\textsf{J}}、$n$ が \abb{\textsf{HI}} の倍数である $n$ の個数\\

\noindent
は \ab{\textsf{K}}、$n + 1$ が \abb{\textsf{HI}} の倍数である $n$ の個数は \ab{\textsf{L}}である。

\vspace{2em}

よって、求める $n$ の個数は \ab{\textsf{MN}} であり、このうち、素数である $n$ は小さい順に \ab{\textsf{OP}} 、\ab{\textsf{QR}} 、\ab{\textsf{ST}} である。

\newpage
\begin{center}
- 計算欄 (memo) -
\end{center}
\newpage
%第四部分
\noindent
\begin{tikzpicture}
\node[draw, rectangle, minimum width=1cm, minimum height=0.8cm, font=\Huge] {IV};
\end{tikzpicture}
\\

2 つの実数 $x, y$ が方程式

\vspace{1em}

$$3x^2 + 2xy + 3y^2 = 32 \hspace{3em} \cdots\cdots \maru{1}$$ 

\vspace{1em}

\noindent
を満たしている。このとき、$x + y$ および $xy$ がとる値の範囲を求めよう。

\vspace{2em}

まず

\vspace{1em}

\hspace{3em} $x + y = a$ \hspace{3em} $\cdots\cdots$ \maru{2}

\vspace{1em}

\noindent
とおく。\maru{1}、\maru{2} より $y$ を消去して、$x$ の 2 次方程式

\vspace{1em}

$$ \ab{\textsf{A}}\ x^2 -\ \ab{\textsf{B}}\ ax +\ \ab{\textsf{C}}\ a^2 - 32 = 0$$

\vspace{1em}

\noindent
を得る。$x$ は実数であるから

\vspace{1em}

$$\ab{\textsf{DE}} \leq a \leq \ab{\textsf{F}} \hspace{3em} \cdots\cdots \maru{3}$$

\vspace{1em}

\noindent
である。

\vspace{1em}

さらに

\vspace{1em}


$$xy = b \hspace{8em} \cdots\cdots \maru{4}$$

\vspace{1em}

\noindent
とおくと、\maru{1}、\maru{2}、\maru{4} より

\vspace{1em}

$$b = \frac{\ab{\textsf{G}}}{\ab{\textsf{H}}}a^2 - \ab{\textsf{I}} \hspace{3em} \cdots\cdots \maru{5}$$

\vspace{1em}

\noindent
を得る。よって、\maru{3}、\maru{5} より

\vspace{1em}

$$\ab{\textsf{JK}} \leq b \leq \ab{\textsf{L}}$$

\vspace{1em}

\noindent
となる
\newpage
\begin{center}
- 計算欄 (memo) -
\end{center}

\end{CJK}

\end{document}